% !TeX program = pdflatex
% Tom IV: Statystyka matematyczna, rozdzialy 1--32 (po audycie 29--30).
% Samodzielny skrypt. Czesc klasyczna i wnioskowanie bayesowskie.
\documentclass[11pt,a4paper,oneside]{report}
\usepackage[T1]{fontenc}
\usepackage[utf8]{inputenc}
\usepackage[polish]{babel}
\usepackage{lmodern}
\usepackage{amsmath,amssymb,amsthm,mathtools}
\usepackage{microtype}
\usepackage{geometry}
\usepackage{enumitem}
\usepackage{booktabs,array,longtable}
\usepackage{xcolor}
\usepackage{needspace}
\usepackage{listings}
\definecolor{codebg}{RGB}{248,249,251}
\definecolor{codeborder}{RGB}{180,191,203}
\definecolor{codecomment}{RGB}{61,113,79}
\definecolor{codekeyword}{RGB}{41,78,132}
\lstdefinelanguage{Julia}{
  morekeywords={using,import,begin,end,if,else,elseif,for,while,return,function,
    struct,mutable,where,let,do,try,catch,finally,const,local,global,println,print},
  sensitive=true,morecomment=[l]{\#},morestring=[b]",
}
\lstdefinestyle{skrypt}{
 backgroundcolor=\color{codebg},frame=single,rulecolor=\color{codeborder},
 framesep=5pt,frameround=tttt,
 basicstyle=\ttfamily\footnotesize,keywordstyle=\color{codekeyword}\bfseries,
 commentstyle=\color{codecomment}\itshape,stringstyle=\color{purple!65!black},
 numbers=left,numberstyle=\tiny\color{gray!75!black},numbersep=8pt,
 stepnumber=1,breaklines=true,breakatwhitespace=true,
 columns=fullflexible,keepspaces=true,showstringspaces=false,
 linewidth=.90\textwidth,
 captionpos=b,aboveskip=11pt,belowskip=12pt,
 xleftmargin=12pt,xrightmargin=2pt,
}
\lstset{style=skrypt}
\usepackage{tikz}
\usetikzlibrary{arrows.meta,positioning,calc,shapes.geometric,fit,backgrounds}
\usepackage[unicode]{hyperref}
% Odsylacze miedzytomowe; zbudowac wszystkie .aux i przechowywac PDF-y obok siebie.
\usepackage{xr-hyper}
\externaldocument[I-]{rachunek_prawdopodobienstwa_tom_I}[rachunek_prawdopodobienstwa_tom_I.pdf]
\externaldocument[II-]{rachunek_prawdopodobienstwa_tom_II}[rachunek_prawdopodobienstwa_tom_II.pdf]
\externaldocument[III-]{rachunek_prawdopodobienstwa_tom_III}[rachunek_prawdopodobienstwa_tom_III.pdf]

\geometry{margin=2.55cm}
\hypersetup{colorlinks=true,linkcolor=blue!45!black,urlcolor=blue!45!black}
\setlength{\parindent}{0pt}
\setlength{\parskip}{4pt}
\setlength{\emergencystretch}{2em}
\setlist[enumerate]{topsep=4pt,itemsep=2pt}
\setlist[itemize]{topsep=4pt,itemsep=2pt}
\numberwithin{equation}{chapter}
\theoremstyle{plain}
\newtheorem{twierdzenie}{Twierdzenie}[chapter]
\newtheorem{lemat}[twierdzenie]{Lemat}
\newtheorem{wniosek}[twierdzenie]{Wniosek}
\newtheorem{stwierdzenie}[twierdzenie]{Stwierdzenie}
\theoremstyle{definition}
\newtheorem{definicja}{Definicja}[chapter]
\newtheorem{przyklad}{Przykład}[chapter]
\newtheorem{zadanie}{Zadanie rozwiązane}[chapter]
\theoremstyle{remark}
\newtheorem{uwaga}{Uwaga}[chapter]
\renewcommand{\proofname}{Dowód}
\newcommand{\R}{\mathbb R}
\newcommand{\N}{\mathbb N}
\newcommand{\Prob}{\mathbb P}
\newcommand{\E}{\mathbb E}
\newcommand{\Var}{\operatorname{Var}}
\newcommand{\Cov}{\operatorname{Cov}}
\newcommand{\Law}{\operatorname{Law}}
\newcommand{\calF}{\mathcal F}
\newcommand{\calA}{\mathcal A}
\newcommand{\calB}{\mathcal B}
\newcommand{\calP}{\mathcal P}
\newcommand{\ind}[1]{\mathbf 1_{#1}}
\newcommand{\dd}{\mathrm d}
\newcommand{\norm}[1]{\lVert#1\rVert}
\newcommand{\abs}[1]{\lvert#1\rvert}
\newcommand{\Normal}{\mathcal N}
\newcommand{\iid}{\mathrel{\overset{\mathrm{iid}}\sim}}
\DeclareMathOperator{\Bin}{Bin}
\DeclareMathOperator{\Pois}{Pois}
\DeclareMathOperator{\Med}{Med}
\DeclareMathOperator{\MSE}{MSE}
\DeclareMathOperator{\supp}{supp}
\DeclareMathOperator{\diag}{diag}
\renewcommand{\lstlistingname}{Kod}
\renewcommand{\lstlistlistingname}{Spis kodów}
\begin{document}
\hypersetup{pageanchor=false}
\begin{titlepage}
\centering\vspace*{3.2cm}
{\Large Podstawy teorii prawdopodobieństwa i statystyki\par}
\vspace{1.1cm}
{\Huge\bfseries Tom IV\\[0.25cm]Statystyka matematyczna\par}
\vspace{0.8cm}
{\LARGE Modele, estymacja, testowanie i decyzje\par}
\vspace{0.7cm}
{\Large Rozdziały 1--32\par}
\vfill
{\large Teoria, dowody, intuicje, dobór metod i przykłady praktyczne\par}
\vspace{1cm}
\end{titlepage}
\hypersetup{pageanchor=true}
\tableofcontents
\input{../rozdzialy/rachunek_prawdopodobienstwa_tom_IV/00-wprowadzenie-do-tomu-iv.tex}
\input{../rozdzialy/rachunek_prawdopodobienstwa_tom_IV/01-eksperyment-statystyczny-i-model-danych.tex}
\input{../rozdzialy/rachunek_prawdopodobienstwa_tom_IV/02-statystyki-i-rozklady-z-proby.tex}
\input{../rozdzialy/rachunek_prawdopodobienstwa_tom_IV/03-funkcja-wiarygodnosci-i-informacja-zawarta-w-danych.tex}
\input{../rozdzialy/rachunek_prawdopodobienstwa_tom_IV/04-statystyczna-teoria-decyzji-od-danych-do-dzialania.tex}
\input{../rozdzialy/rachunek_prawdopodobienstwa_tom_IV/05-konstrukcja-i-porownywanie-estymatorow.tex}
\input{../rozdzialy/rachunek_prawdopodobienstwa_tom_IV/06-dostatecznosc-zupelnosc-i-optymalna-estymacja-nieobciazona.tex}
\input{../rozdzialy/rachunek_prawdopodobienstwa_tom_IV/07-informacja-fishera-efektywnosc-i-granice-dokladnosci-estymacji.tex}
\input{../rozdzialy/rachunek_prawdopodobienstwa_tom_IV/08-asymptotyczna-teoria-estymacji-i-niepewnosc-w-praktyce.tex}
\input{../rozdzialy/rachunek_prawdopodobienstwa_tom_IV/09-testowanie-hipotez-konstrukcja-moc-i-interpretacja.tex}
\input{../rozdzialy/rachunek_prawdopodobienstwa_tom_IV/10-testy-parametryczne-srednie-wariancje-czestosci-i-anova.tex}
\input{../rozdzialy/rachunek_prawdopodobienstwa_tom_IV/11-przedzialy-ufnosci-konstrukcja-pokrycie-i-interpretacja.tex}
\input{../rozdzialy/rachunek_prawdopodobienstwa_tom_IV/12-wielokrotne-testowanie-fwer-fdr-i-decyzje-po-selekcji.tex}
\input{../rozdzialy/rachunek_prawdopodobienstwa_tom_IV/13-regresja-liniowa-estymacja-diagnostyka-i-predykcja.tex}
\input{../rozdzialy/rachunek_prawdopodobienstwa_tom_IV/14-uogolnione-modele-liniowe-regresja-logistyczna-i-poissona.tex}
\input{../rozdzialy/rachunek_prawdopodobienstwa_tom_IV/15-wybor-walidacja-i-regularizacja-modeli-statystycznych.tex}
\input{../rozdzialy/rachunek_prawdopodobienstwa_tom_IV/16-wnioskowanie-z-procesow-zaleznych-autokorelacja-i-modele-szere.tex}
\input{../rozdzialy/rachunek_prawdopodobienstwa_tom_IV/17-podstawy-wnioskowania-bayesowskiego-model-prior-i-posterior.tex}
\input{../rozdzialy/rachunek_prawdopodobienstwa_tom_IV/18-rodziny-sprzezone-rachunek-posterioru-i-prognoza.tex}
\input{../rozdzialy/rachunek_prawdopodobienstwa_tom_IV/19-estymacja-i-predykcja-bayesowska-od-posterioru-do-decyzji.tex}
\input{../rozdzialy/rachunek_prawdopodobienstwa_tom_IV/20-bayesowskie-testowanie-hipotez-iloraz-bayesa-i-paradoksy.tex}
\input{../rozdzialy/rachunek_prawdopodobienstwa_tom_IV/21-modele-hierarchiczne-czesciowe-laczenie-informacji-miedzy-grup.tex}
\input{../rozdzialy/rachunek_prawdopodobienstwa_tom_IV/22-asymptotyka-bayesowska-koncentracja-posterioru-i-twierdzenie-b.tex}
\input{../rozdzialy/rachunek_prawdopodobienstwa_tom_IV/23-metody-monte-carlo-symulacja-blad-i-redukcja-wariancji.tex}
\input{../rozdzialy/rachunek_prawdopodobienstwa_tom_IV/24-metody-mcmc-lancuchy-markowa-estymacja-i-diagnostyka.tex}
\input{../rozdzialy/rachunek_prawdopodobienstwa_tom_IV/25-bootstrap-jackknife-i-testy-permutacyjne-wnioskowanie-przez-po.tex}
\input{../rozdzialy/rachunek_prawdopodobienstwa_tom_IV/26-metody-nieparametryczne-i-odporne-rozklad-rangi-i-obserwacje-o.tex}
\input{../rozdzialy/rachunek_prawdopodobienstwa_tom_IV/27-diagnostyka-rozkladu-i-zalozen-normalnosc-wariancje-i-zaleznos.tex}
\input{../rozdzialy/rachunek_prawdopodobienstwa_tom_IV/28-anova-i-ancova-modele-interakcje-porownania-i-zastosowania.tex}
\input{../rozdzialy/rachunek_prawdopodobienstwa_tom_IV/29-manova-i-mancova-wnioskowanie-o-wielu-wynikach-jednoczesnie.tex}
\input{../rozdzialy/rachunek_prawdopodobienstwa_tom_IV/30-pomiary-powtarzane-sferycznosc-i-modele-mieszane.tex}
\input{../rozdzialy/rachunek_prawdopodobienstwa_tom_IV/31-eksperymenty-a-b-testing-i-wnioskowanie-przyczynowe.tex}
\input{../rozdzialy/rachunek_prawdopodobienstwa_tom_IV/32-laboratorium-statystyczne-wybor-metody-analiza-i-raportowanie.tex}
\end{document}
